%% ============================================================================
%% I. Introduction                               target: ~0.75 page incl. Fig. 1
%% Motivate the problem, state the gap, preview the hypotheses (Sec. IV) and
%% the methodology (Sec. V), and list the contributions.
%% ============================================================================
\section{Introduction}
\label{sec:intro}


Deformable linear objects (DLOs) are objects whose measure along a single dimension, commonly referred to as the object's length, is far greater than the that of the other dimensions. This class includes several classes of items and materials such as ropes, wires, and cables, widely used in industrial, medical, and domestic domains, making DLO manipulation an important step for automation. The Different subsets of DLOs differ on their mechanical properties -- resistance to deformation along and around the different axes, elasticity, and plasticity -- which determine the usefulness of each subset for different purposes. 

Elastic rods are DLOs with relatively high elasticity and resistance to external forces, and low plasticity. This makes elastic rods useful for directional devices that need to be navigated through constrained environments, and resist gravity to maintain their orientation. These include, for example, cameras, thermometers, and brushes. DLO manipulation in these settings is complex. Identifying a DLO using a camera under sub-optimal lighting conditions and in the presence of clutter is a hard task, and, in some cases, becomes impossible due to environment objects.  

In this paper we propose a framework for handling elastic rods by utilizing a contact point of the DLOs with the environment. We hypothesize that given a robotic arm grasping an elastic rod that is in contact with a single point in the environment, i.e., an edge of an object, haptic measurements obtained by force/torque (F/T) or tactile sensors and propreioception data are sufficient for computationally estimating the pose of the DLO's tip. We further hypothesize that a small number of motion primitives can be constructed, and those will allow the manipulation of the DLO and, through that, planning for its tip in a constrained environment with no vision of the tip itself. 

The manuscript is organized as follows. In Section \ref{sec:related} we give a short survey of the literature on DLO estimation and manipulation, with an emphasis on work using F/T and tactile sensors, and motion planning that utilizes objects in the environment. In Section \ref{sec:Preliminaries} we lay out mathematical notations and define terms required for the next sections and formally describe the problem addressed in this paper. Finally, in Sections \ref{sec:hypotheses}, \ref{sec:method}, and \ref{sec:experiments} we give a detailed account of our hypotheses, algorithmic approach, and the experimental evaluation process designed to gather the evidence that will either prove or refute the hypotheses, respectively. 









\begin{outline}
	
  \item \textbf{Motivation.} Cameras, brushes, and other tools are mounted on
    the tips of elastic rods and other DLOs and inserted into constrained
    spaces, \eg for inspection and repair inside engines~\cite{dong2017slender}.
    \todo{1--2 more application examples.}
  \item \textbf{Challenge.} The rod's shape depends on mechanics that are hard
    to model~\cite{sanchez2018survey}, and once inserted, the rod is hidden by
    the environment or by the robot itself, so visual state estimation is
    unreliable.
  \item \textbf{Key idea.} Contact with the environment is usually avoided.
    Instead, we use a contact point between the rod and an environment edge as
    an \emph{ad-hoc joint}: pivoting about it changes the tip direction in ways
    that are impossible if no collisions are tolerated (\cref{fig:teaser}).
  \item \textbf{Problem.} Bring the tip of a gripped elastic rod to a target
    pose in a constrained environment by exploiting a contact with an
    environment edge, using only haptic sensing (robot proprioception, wrist
    force/torque, contact location) during execution (\cref{sec:problem}).
  \item \textbf{Approach.} A sensorized edge reports where the rod touches it;
    a learned estimator maps gripper pose, wrench, and contact location to the
    tip pose; a contact-aware controller steers the tip to the target
    (\cref{sec:method}).
\end{outline}

% Fig. 1 is defined here, not at the top of the section, so that it lands at
% the top of the right column of page 1. LaTeX puts a [t] float at the top of
% the left column (above the abstract) if that column still has room when the
% float is reached; if this happens after edits, move this block further down.
\begin{figure}[t]
  \centering
  \figplaceholder{\linewidth}{4.2cm}{%
    \textbf{Fig.~1 (teaser).} A robot grasps an elastic rod whose tip carries a
    tool. The rod rests on an environment edge (sensorized ring) that acts as an
    ad-hoc joint, letting the tip turn toward a target pose that is unreachable
    without contact. Inset: reachable tip directions with vs.\ without contact.}
  \caption{Using a contact point with an environment edge as an ad-hoc joint
    lets the tip of a gripped elastic rod reach poses that are infeasible
    without contact.}
  \label{fig:teaser}
\end{figure}

\begin{outline}
  \item \textbf{Hypotheses} (\cref{sec:hypotheses}): contact enlarges the set
    of reachable tip poses (\ref{hyp:reach}); haptic signals suffice to estimate
    the tip pose without vision (\ref{hyp:estimate}); closed-loop haptic
    steering reaches targets that contact-free or open-loop strategies miss
    (\ref{hyp:steer}).
  \item \textbf{Contributions}:
    \begin{itemize}
      \item a formulation of DLO tip manipulation that uses a single
        environment contact as an ad-hoc joint, with a reachability analysis;
      \item a vision-free tip-pose estimator that uses proprioception,
        force/torque, and contact location;
      \item an experimental methodology and dataset with a Nitinol rod, from
        planar (2D) to spatial (3D) settings;
      \item preliminary results on the hypotheses.
        \todo{Update after the experiments.}
    \end{itemize}
\end{outline}
